\chapter*{Conclusion}
\addcontentsline{toc}{chapter}{Conclusion}

L'objectif principal de ce mémoire était d'aboutir aux démonstrations du
théorème d'Ax-Grothendieck et du Nullstellensatz. Pour y arriver, nous avons
introduit au fur et à mesure les éléments nécessaires de théorie des modèles.
L'étude des corps algébriquement clos nous a permis de voir comment des notions
de logique comme la compacité, la complétude ou l'élimination des
quantificateurs se traduisent en des outils plus concrets comme le principe de
Lefschetz et permettent par la suite de résoudre des problèmes a priori sans
rapport avec la logique. 

Ce travail montre en particulier que la théorie des modèles s'impose comme un outil géométrique
puissant. Du point de vue de l'algèbre, les résultats que nous avons exposés
constituent le point de départ naturel vers la géométrie algébrique. Le
Nullstellensatz agit comme le dictionnaire fondamental reliant la géométrie à
l'algèbre commutative. Cette correspondance nous invite à étudier les variétés
abstraites via la topologie de Zariski, un contexte abstrait auquel la théorie
des modèles s'adapte d'ailleurs tout naturellement. 

Néanmoins, définir la théorie des modèles comme une boîte à outils de la géométrie algébrique serait restrictif.
Au-delà des applications immédiates, la théorie des modèles pure a développé ses propres dynamiques internes.
Un représentant majeur est la théorie de la classification de Shelah : cette théorie invite à ne pas étudier des théories isolées 
comme nous l'avons fait pour ACF, mais à mettre en lumière des invariants fondamentaux via des propriétés telles que la catégoricité.
Cette classification fondamentale confère à la théorie des modèles une place autonome et une portée universelle au sein des mathématiques.